\chapter{Modelo de negocio y viabilidad econ\'omica}\label{ch:negocio}

\corr{Este cap\'{\i}tulo complementa el an\'alisis de competencias del Cap\'{\i}tulo~\ref{ch:descripcion} (FODA, Porter y Triple~P) con un modelo de sustentabilidad y un an\'alisis econ\'omico-financiero en d\'olares estadounidenses.
Los n\'umeros son supuestos explicitados, coherentes con el alcance del prototipo y con la evidencia de user research; no representan un contrato comercial firmado.}

\section{\corr{Qui\'en usa, qui\'en decide y qui\'en paga}}

\corr{Prevex distingue con claridad tres roles econ\'omicos \parencite{osterwalder2010}:}

\begin{table}[htbp]
	\centering
	\caption[Usuario, decisor y pagador de Prevex]{\corr{Usuario, decisor y pagador de Prevex.} \fuentepropia}
	\label{tab:negocio-roles}
	\small
	\begin{tabular}{@{}p{0.22\textwidth}p{0.68\textwidth}@{}}
		\hline
		\textbf{Rol} & \textbf{Descripci\'on} \\
		\hline
		Usuario & Preventores y coordinadores de la ART: priorizan visitas, preparan intervenciones con el copiloto normativo, registran actas y reciben asignaciones. \\
		Decisor & Gerencia o jefatura de prevenci\'on de la ART: aprueba la adopci\'on si el r\'edito operativo (cobertura, foco y trazabilidad) justifica el gasto. \\
		Pagador & La ART, mediante suscripci\'on B2B. Es quien organiza las visitas, asume obligaciones regulatorias y cubre costos de siniestralidad. \\
		\hline
	\end{tabular}
\end{table}

\corr{La cadena de valor se extiende m\'as all\'a del pagador, sin confundir beneficiario con cliente.
La ART gana eficiencia operativa (ordenamiento de visitas, dossier y normativa a mano, asignaci\'on asistida) y reduce exposici\'on a visitas mal planificadas y a costos de siniestros.
Los establecimientos se benefician si aplican las recomendaciones del preventor (menos ausentismo, mayor productividad).
Los trabajadores se benefician si el entorno laboral mejora y baja la probabilidad de accidentes leves, graves o fatales.
El contrato, no obstante, lo paga la ART.}

\section{\corr{Business Model Canvas}}

\corr{El lienzo de modelo de negocio \parencite{osterwalder2010} resume c\'omo Prevex crea, entrega y captura valor.
Los bloques de segmentos, propuesta de valor, ingresos y costos se apoyan en la encuesta, las entrevistas y el an\'alisis competitivo ya documentados.}

\begin{landscape}
\begin{figure}[htbp]
	\centering
	\corr{\resizebox{0.95\linewidth}{!}{% Business Model Canvas Prevex (layout Osterwalder / Strategyzer)
\begingroup
\definecolor{PrevexBlue}{RGB}{12,74,140}
\definecolor{BMCLine}{RGB}{55,65,81}
\definecolor{BMCHead}{RGB}{15,23,42}
\tikzset{
  bmcbox/.style={
    draw=BMCLine,
    line width=0.7pt,
    rectangle,
    inner sep=5pt,
    align=left,
    font=\sffamily\footnotesize,
    text=BMCHead,
    anchor=north west,
  },
}
% Titulo de bloque: icono azul Prevex + texto en negrita
\newcommand{\bmchead}[2]{%
  {\color{PrevexBlue}\Large#1}\hspace{0.35em}%
  {\sffamily\bfseries\small\color{PrevexBlue}#2}%
}
\begin{tikzpicture}[x=1cm, y=1cm]
  \def\Wsoc{4.6}
  \def\Wmid{4.6}
  \def\Wval{5.8}
  \def\Htop{7.2}
  \def\Hhalf{3.52}
  \def\Hbot{3.1}
  \def\gap{0.12}

  \def\xsoc{0}
  \pgfmathsetmacro{\xact}{\Wsoc+\gap}
  \pgfmathsetmacro{\xval}{\xact+\Wmid+\gap}
  \pgfmathsetmacro{\xrel}{\xval+\Wval+\gap}
  \pgfmathsetmacro{\xseg}{\xrel+\Wmid+\gap}
  \pgfmathsetmacro{\Wtot}{\xseg+\Wsoc}
  \pgfmathsetmacro{\Wcost}{(\Wtot-\gap)/2}
  \pgfmathsetmacro{\xrev}{\Wcost+\gap}

  \node[anchor=south west, font=\sffamily\bfseries\large\color{PrevexBlue}, inner sep=0] at (0,7.85)
    {Business Model Canvas: Prevex};

  \node[bmcbox, minimum width=\Wsoc cm, minimum height=\Htop cm,
        text width=\dimexpr\Wsoc cm-12pt\relax] (soc) at (\xsoc,7.65) {%
    \bmchead{\faHandshake}{Socios clave}\\[5pt]
    ART cliente; integradores de core ART; proveedor cloud; API de modelo generativo.
  };

  \node[bmcbox, minimum width=\Wmid cm, minimum height=\Hhalf cm,
        text width=\dimexpr\Wmid cm-12pt\relax] (act) at (\xact,7.65) {%
    \bmchead{\faTasks}{Actividades clave}\\[5pt]
    Pipeline de datos; entrenamiento XGBoost; operaci\'on API/frontend; copiloto normativo y ranking.
  };

  \node[bmcbox, minimum width=\Wmid cm, minimum height=\Hhalf cm,
        text width=\dimexpr\Wmid cm-12pt\relax] (rec) at (\xact,7.65-\Hhalf-\gap) {%
    \bmchead{\faCogs}{Recursos clave}\\[5pt]
    Panel tabular; modelo y SHAP; corpus SRT; infraestructura cloud; soporte.
  };

  \node[bmcbox, minimum width=\Wval cm, minimum height=\Htop cm,
        text width=\dimexpr\Wval cm-12pt\relax] (val) at (\xval,7.65) {%
    \bmchead{\faGift}{Propuestas de valor}\\[5pt]
    Ranking trimestral explicable; preparaci\'on de visita con normativa SRT y dossier; cierre estructurado; asignaci\'on asistida. Apoyo a la decisi\'on, sin reemplazar al preventor.
  };

  \node[bmcbox, minimum width=\Wmid cm, minimum height=\Hhalf cm,
        text width=\dimexpr\Wmid cm-12pt\relax] (rel) at (\xrel,7.65) {%
    \bmchead{\faHeart}{Relaci\'on con clientes}\\[5pt]
    Onboarding; soporte operativo; reentrenamiento del modelo; actualizaci\'on del corpus normativo.
  };

  \node[bmcbox, minimum width=\Wmid cm, minimum height=\Hhalf cm,
        text width=\dimexpr\Wmid cm-12pt\relax] (can) at (\xrel,7.65-\Hhalf-\gap) {%
    \bmchead{\faTruck}{Canales}\\[5pt]
    Venta B2B a \'areas de prevenci\'on; canal con proveedores de core ART.
  };

  \node[bmcbox, minimum width=\Wsoc cm, minimum height=\Htop cm,
        text width=\dimexpr\Wsoc cm-12pt\relax] (seg) at (\xseg,7.65) {%
    \bmchead{\faUsers}{Segmentos de clientes}\\[5pt]
    ART en Argentina con cartera de establecimientos y equipo de prevenci\'on (preventores y coordinadores). Pagador: la ART.
  };

  \pgfmathsetmacro{\ybot}{7.65-\Htop-\gap}
  \node[bmcbox, minimum width=\Wcost cm, minimum height=\Hbot cm,
        text width=\dimexpr\Wcost cm-12pt\relax] (cost) at (0,\ybot) {%
    \bmchead{\faTags}{Estructura de costos}\\[5pt]
    Fijos: cloud base, soporte, reentrenamiento. Variables: inferencia generativa, almacenamiento, CAC por ART nueva. Inversi\'on inicial de productizaci\'on.
  };

  \node[bmcbox, minimum width=\Wcost cm, minimum height=\Hbot cm,
        text width=\dimexpr\Wcost cm-12pt\relax] (ing) at (\xrev,\ybot) {%
    \bmchead{\faMoneyBillWave}{Fuentes de ingresos}\\[5pt]
    Suscripci\'on SaaS mensual o anual por ART (un contrato por aseguradora), no por preventor. Precio de referencia base: USD~1.200/mes.
  };

  \node[anchor=north west, font=\sffamily\scriptsize\color{BMCLine}, inner sep=0]
    at (0,\ybot-\Hbot-0.22)
    {Fuente: adaptado de Osterwalder y Pigneur, Business Model Generation (2010).};
\end{tikzpicture}
\endgroup
}}
	\caption[Business Model Canvas de Prevex]{\corr{Business Model Canvas de Prevex.} \fuentepropia}
	\label{fig:negocio-canvas}
\end{figure}
\end{landscape}

\section{\corr{Fuentes de ingresos y estructura de costos}}

\corr{La unidad de venta es la ART.
El mercado de referencia son las \textbf{22 aseguradoras registradas} publicadas por el Estado nacional para elecci\'on del empleador \parencite{srtartlistado}.
El precio de referencia es \textbf{el mismo} en los tres escenarios: USD~500 por mes (USD~6.000 por a\~no).
Ese precio se combina con una estructura de costos que ya incluye sueldos (no solo cloud).
Lo que diferencia pesimista, base y optimista es la \textbf{cartera fija} de clientes (3, 5 o 7 ART durante los cinco a\~nos).}

\corr{Los costos se clasifican as\'{\i}:}
\begin{itemize}
	\item \corr{\textbf{Inversi\'on inicial (capex):} USD~8.000 de productizaci\'on m\'{\i}nima (despliegue endurecido, aspectos legales livianos, materiales comerciales y primer onboarding). No es solo la factura de cloud ni de la API generativa, y no incluye las horas del PFI.}
	\item \corr{\textbf{Costos fijos anuales (opex fijo):} USD~15.000 en total, compuestos por infraestructura cloud, monitoreo y reentrenamiento (USD~3.000) m\'as \textbf{sueldos} de un perfil part-time de operaciones y soporte (USD~12.000). Sin esa partida el modelo subestimaba el costo real de sostener el producto.}
	\item \corr{\textbf{Costos variables por ART:} uso de API generativa y almacenamiento incremental (USD~800 por ART-a\~no).}
	\item \corr{\textbf{CAC:} costo de adquisici\'on al incorporar una ART nueva (USD~3.000). En cartera fija se concentra en el a\~no~1.}
\end{itemize}

\begin{table}[htbp]
\centering
\caption[Supuestos de los escenarios (USD)]{\rev{Supuestos de los escenarios (USD).} \fuentepropia}
\label{tab:negocio-supuestos}
\small
\begin{tabular}{lrrr}
\hline
\textbf{Supuesto} & \textbf{Pesimista} & \textbf{Base} & \textbf{Optimista} \\
\hline
Precio mensual / ART & USD~500 & USD~500 & USD~500 \\
Opex fijo anual & USD~3.000 & USD~3.000 & USD~3.000 \\
Costo variable / ART-ano & USD~800 & USD~800 & USD~800 \\
CAC por ART nueva & USD~3.000 & USD~3.000 & USD~3.000 \\
ART fijas (5 a\~nos) & 1 & 3 & 5 \\
\hline
\end{tabular}
\end{table}


\section{\corr{Viabilidad econ\'omico-financiera}}

\corr{El an\'alisis se realiza en USD (moneda estable), con horizonte de cinco a\~nos y tasa de descuento $k=18\%$ real.
Se reportan VAN, TIR, payback simple y descontado, y punto de equilibrio en n\'umero de ART.}

\corr{Se construyen tres escenarios con el mismo precio y distinta cartera fija (fracci\'on del mercado de~22 ART):}

\begin{itemize}
	\item \corr{\textbf{Pesimista:} tres ART fijas (~14\,\% del mercado) a USD~500/mes.}
	\item \corr{\textbf{Base:} cinco ART fijas (~23\,\% del mercado), mismo precio.}
	\item \corr{\textbf{Optimista:} siete ART fijas (~32\,\% del mercado), mismo precio.}
\end{itemize}

\begin{table}[htbp]
\centering
\caption[Flujo de fondos del escenario base (USD)]{\rev{Flujo de fondos del escenario base (USD). A\~no 0 = inversi\'on inicial.} \fuentepropia}
\label{tab:negocio-flujo-base}
\small
\begin{tabular}{rrrrrr}
\hline
\textbf{A\~no} & \textbf{ART} & \textbf{Ingresos} & \textbf{Costos} & \textbf{Flujo neto} & \textbf{Acumulado} \\
\hline
0 & 0 & USD~0 & USD~8.000 & -USD~8.000 & -USD~8.000 \\
1 & 1 & USD~14.400 & USD~14.800 & -USD~400 & -USD~8.400 \\
2 & 2 & USD~28.800 & USD~16.100 & USD~12.700 & USD~4.300 \\
3 & 3 & USD~43.200 & USD~17.400 & USD~25.800 & USD~30.100 \\
4 & 4 & USD~57.600 & USD~18.700 & USD~38.900 & USD~69.000 \\
5 & 5 & USD~72.000 & USD~20.000 & USD~52.000 & USD~121.000 \\
\hline
\end{tabular}
\end{table}


\begin{table}[htbp]
\centering
\caption[Indicadores de viabilidad por escenario (USD)]{\rev{Indicadores de viabilidad por escenario (USD, $k=12\%$, horizonte 5 a\~nos).} \fuentepropia}
\label{tab:negocio-indicadores}
\small
\begin{tabular}{lrrr}
\hline
\textbf{Indicador} & \textbf{Pesimista} & \textbf{Base} & \textbf{Optimista} \\
\hline
VAN & -USD~12.395 & USD~74.359 & USD~259.174 \\
TIR & n/d & 115.9\% & 318.5\% \\
Payback simple (a\~nos) & no recupera en 5 a\~nos & 1.7 & 0.5 \\
Payback descontado (a\~nos) & no recupera en 5 a\~nos & 1.8 & 0.6 \\
Punto de equilibrio (ART) & 1.1 & 0.8 & 0.7 \\
\hline
\end{tabular}
\end{table}


\corr{Interpretaci\'on: con sueldos incluidos, el punto de equilibrio ronda \textbf{2,9~ART}.
El pesimista (tres ART) no cierra: VAN de USD~-13.751 y no recupera la inversi\'on en cinco a\~nos.
En el escenario base el VAN es positivo (USD~13.687), la TIR (~55\,\%) supera a $k$ y el payback descontado ronda 2,5~a\~nos.
El optimista, con siete clientes, refuerza la lectura (VAN de USD~41.125; payback descontado de~1,5~a\~nos).
La lectura comercial es que el portafolio chico por debajo del equilibrio no alcanza cuando se pagan sueldos; con cinco ART el proyecto ya genera valor.}

\corr{Una sensibilidad de $k$ entre~15\% y~20\% no cambia el signo del VAN en base ni en optimista; el pesimista permanece negativo.
Los n\'umeros siguen siendo supuestos expl\'{\i}citos, no un contrato firmado.}

\section{\corr{Branding e identidad visual}}

\corr{El nombre \emph{Prevex} combina prevenci\'on con explicabilidad: el producto no solo predice, sino que expone los factores que ordenan la lista (SHAP) y apoya la visita con evidencia normativa.}

\begin{figure}[htbp]
	\centering
	\includegraphics[width=0.42\textwidth]{figures/logo_prevex.png}
	\caption[Logotipo de Prevex]{\corr{Logotipo de Prevex utilizado en la interfaz del prototipo.} \fuentepropia}
	\label{fig:negocio-logo}
\end{figure}

\corr{La marca aparece en login, encabezado y materiales de la aplicaci\'on; en este documento se presenta como identidad del prototipo evaluado, no como campa\~na comercial.}
